\documentclass[10pt,a4paper]{amsart}
\usepackage[margin=19mm]{geometry}
\usepackage{amsmath,amssymb,mathtools}
\usepackage[scr=rsfs,cal=euler]{mathalfa}
\usepackage{xfrac}
\usepackage[unicode,hidelinks,bookmarks=false]{hyperref}
\newcommand{\ie}{i.e.\ }
\newcommand{\cf}{cf.\ }
\newcommand{\resp}{respectively\ }
\newcommand{\Subsection}{\subsection}
\providecommand{\nobreakdash}{\nobreak\hbox{-}}
\setlength{\emergencystretch}{2em}
\multlinegap0pt
\usepackage{bm}
\usepackage{bbm}
\usepackage{dsfont}
\usepackage{xspace}
\usepackage{cancel}
\usepackage{mathtools}
\usepackage{amsthm}
\makeatletter
\@ifundefined{theorem}{\newtheorem{theorem}{Theorem}}{}
\@ifundefined{proposition}{\newtheorem{proposition}[theorem]{Proposition}}{}
\makeatother
\title[Rook theory, normal ordering in the $q$-deformed Ore algebra]{Rook theory, normal ordering in the $q$-deformed Ore algebra and the polynomial generalization}
\author{Matthias Schork}
\date{Source: arXiv:2605.09683v3 (CC BY 4.0)}
\begin{document}
\maketitle

\noindent\emph{Source and licence.} Rook theory, normal ordering in the $q$-deformed Ore algebra and the polynomial generalization. Version arXiv:2605.09683v3 (CC BY 4.0), distributed under the Creative Commons Attribution 4.0 International licence, as recorded in the arXiv licence metadata for this exact version and shown on the abstract page https://arxiv.org/abs/2605.09683v3. Full rights notice: licenses/arxiv-2605.09683-CC-BY-4.0.txt.\par\smallskip

\noindent\textit{Editorial scope.} Counted unit: the Proposition whose source label is \texttt{THRECPOLY} in the article (source0.tex lines 1266--1270, source label \texttt{THRECPOLY}), delivered together with the article's own complete original proof of it (source0.tex lines 1271--1303, one \texttt{proof} environment of 33 lines). No mathematics is written, repaired or re-derived: statement and proof are the article's own text line for line. Required context retained uncounted: PROPQORESRec (lines 726--730); PolyStir (lines 1241--1245). Every definition, display and label of those lines is retained unchanged. Omitted from this package and not redistributed: 1--725, 731--1240, 1246--1265, 1304--1545. Nothing inside a retained excerpt is omitted or altered: every retained body is delivered exactly as the article's lines are written, so no omission receipt of any kind accompanies this package. Citations: the retained excerpts carry no citation command at all, so no bibliography entry of the article is transcribed and no reference is redistributed. Reference adaptation: none was needed and none was made; every reference inside the delivered unit resolves inside it, so no unresolved reference or citation is produced. Printed slips kept literal: no printed word, symbol, hypothesis or number of the counted statement or of its proof was altered. Layout and class: the article's own class is \texttt{amsart} with packages including amssymb, latexsym, pstricks, color, hyperref, tensor, enumitem, mathabx, comment, amsfonts, amsmath, amsthm, graphicx, mathrsfs; the wrapper is the lane's pinned \texttt{amsart} editorial wrapper at 10pt on A4, which restates the theorem-like environments the retained text needs and adds the article's own macro definitions that the retained text needs (none), with extra wrapper packages (bm, bbm, dsfont, xspace, cancel, mathtools). The delivered numbering is the unit's own. Assets: no retained span contains a figure environment, a graphics-inclusion command, a PDF-page inclusion, a table or a TikZ picture; no image file is copied into this package, the package declares no asset dependency and no image file is redistributed with it. Licence: version arXiv:2605.09683v3 of this article is distributed under the Creative Commons Attribution 4.0 International licence, as recorded in the arXiv licence metadata for this exact version and shown on the abstract page https://arxiv.org/abs/2605.09683v3. A full rights notice is delivered at licenses/arxiv-2605.09683-CC-BY-4.0.txt. Characters: every retained excerpt is pure ASCII, so the delivered engine, XeLaTeX with its default Latin Modern text font, reports no missing character.
\begin{proposition}\label{PROPQORESRec} The $q$-deformed Ore-Stirling numbers satisfy the recurrence relation
\begin{equation}\label{QORESRec}
S_{\mu,\nu;q}(n+1;j,k) =q^{j-1}S_{\mu,\nu;q}(n;j-1,k-1)+\mu [j]_q S_{\mu,\nu;q}(n;j,k) +\nu [j-1]_q S_{\mu,\nu;q}(n;j-1,k).  
\end{equation}
\end{proposition}

\begin{proposition} The $q$-deformed polynomial Stirling numbers of degree $s$ are given by 
\begin{equation}\label{PolyStir}
{\mathscr S}_{{\bf \alpha};q}(n;j,k)=\sum_{{\bf k} \in {\mathcal WC}_{s+1}(n-k|j-n)} m_{{\bf k}}(J_n;q).
\end{equation}
\end{proposition}

\begin{proposition}\label{THRECPOLY} The $q$-deformed polynomial Stirling numbers of degree $s$ satisfy the recurrence relation 
\begin{equation}\label{RECPOLY}
{\mathscr S}_{{\bf \alpha};q}(n+1;j,k)=q^{j-1}{\mathscr S}_{{\bf \alpha};q}(n;j-1,k-1)+\sum_{r=0}^s \alpha_r [j-r]_q {\mathscr S}_{{\bf \alpha};q}(n;j-r,k).
\end{equation}
\end{proposition}

\begin{proof}
The proof is combinatorial and very similar to the one of Proposition~\ref{PROPQORESRec}. Using \eqref{PolyStir} and the definition of the $m_{{\bf k}}(B;q)$, one has
$$
{\mathscr S}_{{\bf \alpha};q}(n+1;j,k)=\sum_{{\bf k} \in {\mathcal WC}_{s+1}(n+1-k|j-n-1)}{\bf  \alpha}^{{\bf k}} \sum_{\phi \in {\mathcal M}_{{\bf k}}(J_{n+1})} q^{\# \mbox{empty boxes}}.
$$ 
Thus, we have to consider all mixed placements $\phi \in \mathcal{M}_{{\bf k}}(J_{n+1})$. As above, we write in an informal fashion $J_{n+1}=C_n \oplus J_n$. For each type ${\bf k}$, there are $(s+2)$ types of placements of the rooks of the different weights on $J_{n+1}$:
\begin{itemize}
\item Type $I$: Place all rooks on $J_n$ and leave $C_n$ empty.
\item Type $II_r$: Place one rook of weight $r$ in $C_n$ and the remaining rooks on $J_{n}$, for $r=0,\ldots,s$.
\end{itemize} 
Thus, we can write
$$
\mathcal{M}_{{\bf k}}(J_{n+1})=\mathcal{M}_{{\bf k}}^{I}(J_{n+1})\cup\mathcal{M}_{{\bf k}}^{II_0}(J_{n+1})\cup\mathcal{M}_{{\bf k}}^{II_1}(J_{n+1})\cup \cdots \cup\mathcal{M}_{{\bf k}}^{II_s}(J_{n+1}).
$$
Let us start with placements of type $I$. Each such placement  $\phi  \in \mathcal{M}_{{\bf k}}^{I}(J_{n+1})$ corresponds to a placement $\phi' \in \mathcal{M}_{{\bf k}}(J_{n})$. In $C_n$, one has for a placement of type ${\bf k}=(k_0,\ldots,k_s)$ in total $n+\sum_{j=0}^s(j-1)k_j$ empty boxes. Here ${\bf k}\in {\mathcal WC}_{s+1}(n+1-k|j-n-1)$, i.e., $\sum_{j=0}^s(j-1)k_j=(j-n-1)$. Hence, there are $n+(j-n-1)=j-1$ empty boxes in $C_n$. If $\phi'$ has weight $\omega_{{\bf \alpha};q}(J_n,\phi')$ then $\omega_{{\bf \alpha};q}(J_{n+1},\phi)=q^{j-1}\omega_{{\bf \alpha}}(J_n,\phi')$. Thus, the contribution of all placements of type $I$ is given by
$$
\sum_{{\bf k} \in {\mathcal WC}_{s+1}(n-(k-1)|(j-1)-n)}{\bf  \alpha}^{{\bf k}} q^{j-1}\sum_{\phi' \in {\mathcal M}_{{\bf k}}(J_{n})} q^{\# \mbox{empty boxes}} =q^{j-1}{\mathscr S}_{{\bf \alpha};q}(n;j-1,k-1).
$$
Let us turn to placements of type $II_r$ (for $r=0,\ldots,s$). For ${\bf k}=(k_0,\ldots,k_s)$, let us denote 
$$
{\bf k}^{(r)}\equiv (k_0^{(r)},\ldots,k_{s}^{(r)})=(k_0-\delta_{r,0},\ldots,k_{s}-\delta_{r,s}), 
$$
i.e., the $r$-th entry is reduced by one. One has $\sum_{j=0}^s(j-1)k_{j}^{(r)}=\sum_{j=0}^s(j-1)k_{j}-(r-1)$. Thus, 
$$
{\bf k}\in {\mathcal WC}_{s+1}(n+1-k|j-n-1) \Longrightarrow {\bf k}^{(r)}\in {\mathcal WC}_{s+1}(n-k|j-n-r),
$$
and one has ${\bf  \alpha}^{{\bf k}}=\alpha_r{\bf  \alpha}^{{\bf k}^{(r)}}$. If $\phi' \in \mathcal{M}_{{\bf k}^{(r)}}(J_{n})$, then there are $n+\sum_{j=0}^s(j-1)k_{j}^{(r)}=j-r$ possible boxes to place the rook of weight $r$ on $C_n$ to get a placement $\phi \in \mathcal{M}_{{\bf k}}^{II_r}(J_{n+1})$. Placing the rook in the $\ell$-th possible row from above (where $\ell=1,\ldots,j-r$) will result in $(j-r-\ell)$ additional empty boxes. Thus, the sum of the weights of these $j-r$ placements $\phi$ is given by $\alpha_r (1+q+\cdots+q^{j-r-1})\omega_{{\bf \alpha};q}(J_n,\phi')=\alpha_r [j-r]_q \omega_{{\bf \alpha};q}(J_n,\phi')$. The contribution of all placements of type $II_r$ is, therefore, given by
$$
\sum_{{\bf k}^{(r)} \in {\mathcal WC}_{s+1}(n-k|j-r-n)} \alpha_r {\bf  \alpha}^{{\bf k}^{(r)}}[j-r]_q 
\sum_{\phi' \in {\mathcal M}_{{\bf k}^{(r)}}(J_{n})} q^{\# \mbox{empty boxes}}=\alpha_r [j-r]_q {\mathscr S}_{{\bf \alpha};q}(n;j-r,k).
$$
Summing over the types $I$ and $II_r$, for $r=0,\ldots,s$, yields the assertion.
\end{proof}

\end{document}
